\chapter{Remotes and Collaboration}\label{ch:git-remotes}

Everything so far worked with no network. Git is \vocab{distributed}: your clone
is a complete repository with the full history, and a remote is just another
clone you have agreed to exchange commits with.

\section{Remotes}\label{sec:git-remote}

\begin{definition}[Remote]
	A named URL of another repository. \cmd{origin} is the conventional name for
	the one you cloned from; \cmd{upstream} is the convention for the original
	project when you are working from a fork.
\end{definition}

\begin{keys}[0.36]
	\cmd{git remote -v}                   & List remotes and their URLs             \\
	\cmd{git remote add upstream <url>}   & Add one                                 \\
	\cmd{git remote remove upstream}      & Remove one                              \\
	\cmd{git remote show origin}          & What it has, and how your branches track it \\
	\cmd{git remote set-url origin <url>} & Change where it points                  \\
\end{keys}

\begin{shell}
$ git remote -v
origin    git@github.com:Konyi0613/LaTeX.git (fetch)
origin    git@github.com:Konyi0613/LaTeX.git (push)
\end{shell}

\begin{definition}[Remote-tracking branch]
	A local, read-only record of where a branch stood on the remote when you last
	fetched --- written \cmd{origin/main}. It is not the remote's branch, and it
	does not update by itself: it is a cached snapshot.
\end{definition}

\begin{note}
	The distinction explains \emph{``Your branch is ahead of 'origin/main' by 2
		commits''}. Git is comparing your \cmd{main} against its own cached
	\cmd{origin/main}, not against the server. If someone else pushed five minutes
	ago, Git does not know until you fetch.
\end{note}

\section{\texorpdfstring{\cmd{fetch} versus \cmd{pull}}{fetch versus pull}}\label{sec:git-fetch}

\begin{keys}[0.34]
	\cmd{git fetch}                & Download new commits; change nothing else     \\
	\cmd{git fetch -{}-all}        & From every remote                             \\
	\cmd{git fetch -{}-prune}      & \dots{} and delete remote-tracking branches that are gone \\
	\cmd{git pull}                 & \cmd{fetch}, then \cmd{merge}                 \\
	\cmd{git pull -{}-rebase}      & \cmd{fetch}, then \cmd{rebase}                \\
\end{keys}

\begin{moral}
	\cmd{git fetch} is always safe. It updates your remote-tracking branches and
	touches neither your working tree nor your branch, so you can look before
	deciding. \cmd{git pull} is \cmd{fetch} plus an immediate merge, which is why
	it is the command that occasionally produces a conflict you were not expecting.
\end{moral}

\begin{shell}
$ git fetch
$ git log --oneline HEAD..origin/main     # what did they do?
$ git diff HEAD origin/main               # and what does it change?
$ git merge origin/main                   # now decide
\end{shell}

\section{\texorpdfstring{\cmd{push} and Tracking}{push and Tracking}}\label{sec:git-push}

\begin{keys}[0.38]
	\cmd{git push}                          & Push the current branch to its upstream \\
	\cmd{git push -u origin feature}        & Push and set the upstream for next time \\
	\cmd{git push origin -{}-delete feature} & Delete a branch on the remote          \\
	\cmd{git push -{}-tags}                 & Push tags, which \cmd{push} does not do \\
	\cmd{git push -{}-force-with-lease}     & Overwrite the remote branch, but only if nobody else has pushed \\
	\cmd{git push -{}-force}                & Overwrite unconditionally               \\
\end{keys}

\begin{note}[Which force to use]
	After an interactive rebase your local history no longer matches the remote and
	a plain push is rejected. \cmd{--force} overwrites whatever is there, including
	a colleague's commits pushed in the meantime. \cmd{--force-with-lease} refuses
	unless the remote is still where you last saw it --- the same operation, with a
	check. There is no reason to use plain \cmd{--force} on a shared branch.
\end{note}

\begin{note}
	\texttt{fatal: The current branch has no upstream branch} means Git does not
	know where to push. \cmd{git push -u origin feature} sets it once; thereafter
	plain \cmd{git push} works. Setting
	\cmd{git config --global push.autoSetupRemote true} makes Git do it for you.
\end{note}

\section{Merge or Rebase on Pull}\label{sec:git-pullstrat}

Default \cmd{git pull} merges, and on a shared branch this litters history with
merge commits saying nothing but ``I synced'':

\begin{shell}
Merge branch 'main' of github.com:Konyi0613/LaTeX
Merge branch 'main' of github.com:Konyi0613/LaTeX
Merge branch 'main' of github.com:Konyi0613/LaTeX
\end{shell}

\begin{shell}
$ git config --global pull.rebase true       # rebase instead of merging
$ git config --global rebase.autoStash true  # stash and restore automatically
\end{shell}

\begin{moral}
	Rebase when pulling, merge when integrating. Pulling other people's commits
	onto your unpushed work is exactly the case the golden rule of
	\autoref{sec:git-rebase} permits --- your commits are not published yet.
	Bringing a finished branch into \cmd{main} is the case for a merge commit,
	because the fact that it was a branch is worth recording.
\end{moral}

\section{Tags and Releases}\label{sec:git-tags}

\begin{definition}[Tag]
	A name permanently attached to one commit. Unlike a branch, it does not move.
	A \vocab{lightweight} tag is just a ref; an \vocab{annotated} tag is a real
	object with an author, a date and a message, and is what releases should use.
\end{definition}

\begin{keys}[0.36]
	\cmd{git tag}                        & List tags                                \\
	\cmd{git tag -a v1.0 -m "Release 1.0"} & Create an annotated tag                \\
	\cmd{git tag v1.0-beta}              & A lightweight tag                        \\
	\cmd{git tag -a v0.9 a1b2c3d}        & Tag an older commit                      \\
	\cmd{git show v1.0}                  & The tag and the commit it names          \\
	\cmd{git push origin v1.0}           & Push one tag                             \\
	\cmd{git push -{}-follow-tags}       & Push commits and their annotated tags    \\
	\cmd{git tag -d v1.0}                & Delete locally                           \\
	\cmd{git describe -{}-tags}          & ``v1.0-14-ga1b2c3d'': the nearest tag, plus distance \\
\end{keys}

\begin{note}
	\cmd{git describe} is how build systems stamp a version into a binary without
	anyone maintaining a version number by hand. It reads as: fourteen commits
	after \texttt{v1.0}, at commit \texttt{a1b2c3d}.
\end{note}

\section{Submodules and Subtrees}\label{sec:git-sub}

Two ways to put one repository inside another.

\begin{keys}[0.38]
	\cmd{git submodule add <url> lib/foo} & Add a submodule                          \\
	\cmd{git clone -{}-recurse-submodules <url>} & Clone, fetching them too          \\
	\cmd{git submodule update -{}-init -{}-recursive} & Populate them after a plain clone \\
	\cmd{git submodule update -{}-remote}  & Move them to their latest commits       \\
	\cmd{git subtree add -{}-prefix=lib/foo <url> main -{}-squash} & Add a subtree    \\
\end{keys}

\begin{intuition}
	A \vocab{submodule} records a \emph{pointer}: the parent repository stores only
	the URL and one commit hash, and the files live in a separate repository. A
	\vocab{subtree} copies the other project's files into yours, so a plain clone
	gets everything.
\end{intuition}

\begin{note}
	Submodules have a deserved reputation. A clone that omits
	\cmd{--recurse} leaves empty directories; a commit inside a
	submodule is invisible until the parent's pointer is also committed; and
	switching branches routinely leaves submodules on the wrong commit. They are
	right when the dependency is genuinely a separate project with its own
	history; for vendoring a library you do not develop, a subtree or a package
	manager causes less trouble.
\end{note}

\begin{moral}
	Before reaching for either, ask whether your language's package manager solves
	it. Submodules are a version-control answer to a dependency-management
	question, and they are usually the wrong shape for it.
\end{moral}
